\documentclass[11pt,letterpaper]{article}
\usepackage[margin=1in]{geometry}
\usepackage{booktabs}
\usepackage{longtable}
\usepackage{amsmath}
\usepackage{hyperref}
\usepackage{xcolor}
\usepackage{fancyhdr}

\pagestyle{fancy}
\fancyhead[L]{\textsc{DSF-AI Structural Analysis Service}}
\fancyhead[R]{\textsc{Trade Secret --- Do Not Distribute}}
\fancyfoot[C]{\thepage}

\definecolor{confidential}{rgb}{0.7,0,0}

\title{%
  \textbf{DSF-AI Universal Structural Analysis Service} \\[6pt]
  \large Service Specification v1.0 \\[4pt]
  \textcolor{confidential}{\textbf{TRADE SECRET --- CONFIDENTIAL}}
}
\author{Joseph Forrester}
\date{May 4, 2026}

\begin{document}
\maketitle
\thispagestyle{fancy}

\begin{center}
\fbox{\parbox{0.85\textwidth}{\centering
\textbf{Notice:} This document describes a proprietary service built on the
UF-Core kernel (L0--L4) and protected by the Secure Envelope System (SES).
The kernel implementation, intermediate layer outputs, and DSF 7-tuple
geometry are trade secrets. Nothing in this document exposes the kernel
internals.}}
\end{center}

\tableofcontents
\newpage

% ==================================================================
\section{Executive Summary}
% ==================================================================

DSF-AI is a universal structural analysis engine that detects phase
transitions, structural precursors, and regime changes in any
measurement-vs-stimulus data. It operates with zero domain knowledge ---
the same kernel that analyzes financial time series, infrared sensor
calibration curves, and superconductor transport data without modification.

The service accepts a CSV (stimulus column, measurement column) and
returns a structured report identifying every transition, its precursor
onset, the uncertainty profile, and the regime map.

\subsection{Capability Proof}

On May 4, 2026, the kernel was run on six phase transitions across
five materials spanning four distinct physical phenomena and two
measurement types:

\begin{table}[h]
\centering
\begin{tabular}{@{}llll@{}}
\toprule
\textbf{Material} & \textbf{Transition} & \textbf{Type} & \textbf{Detected?} \\
\midrule
YBa$_2$Cu$_3$O$_{7-\delta}$ & $T_c = 93$\,K & High-$T_c$ superconductor & Yes \\
MgB$_2$ & $T_c = 39$\,K & Conventional superconductor & Yes \\
CsV$_3$Sb$_5$ & $T_{c,\text{SC}} = 2.5$\,K & Kagome superconductor & Yes \\
CsV$_3$Sb$_5$ & $T_{\text{CDW}} = 94$\,K & Charge density wave & Subtle \\
VO$_2$ & $T_{\text{MIT}} = 340$\,K & Metal--insulator transition & Yes \\
BaTiO$_3$ & $T_{\text{Curie}} = 393$\,K & Ferroelectric (dielectric) & Yes \\
FeSe/STO & $T_c = 65$\,K & Iron-based superconductor & Yes \\
\bottomrule
\end{tabular}
\caption{Phase transition detection survey. All transitions found from
raw data with zero physics assumptions.}
\end{table}

In every case the kernel also identified \textbf{structural precursors}
above/before the known transition temperature --- signatures that
standard analysis methods either miss or require specialized probes
(ARPES, NMR, tunneling spectroscopy) to detect.

% ==================================================================
\section{What the Service Does}
% ==================================================================

\subsection{Input}

A two-column CSV file:
\begin{itemize}
  \item \textbf{Column 1:} Stimulus (temperature, pressure, voltage,
    time, magnetic field, cycle number, etc.)
  \item \textbf{Column 2:} Measurement (resistance, dielectric constant,
    magnetization, heat flow, voltage, impedance, etc.)
\end{itemize}

No metadata, no material identity, no domain labels required. The kernel
is domain-agnostic.

\subsection{Output}

The service produces two outputs: a structured JSON payload for
programmatic use, and a plain-language narrative summary generated
by an LLM interpretation layer.

\subsubsection{Structured Report (JSON)}

\begin{enumerate}
  \item \textbf{Transition list.} Every structural transition found,
    ranked by severity, with:
    \begin{itemize}
      \item Transition stimulus value (e.g., temperature)
      \item Transition type (regime label: STABLE $\to$ VOLATILE, etc.)
      \item Severity score (coupling strength at transition)
    \end{itemize}
  \item \textbf{Precursor onset.} For each transition, the stimulus
    value where the structural state \emph{begins} changing --- before
    the transition itself. This is the service's primary novel capability.
  \item \textbf{Uncertainty profile.} Per-gate $U_k$ values showing
    where the data is structurally ambiguous.
  \item \textbf{Regime map.} Stimulus ranges classified as STABLE,
    VOLATILE, TRANSITIONAL, or DEGENERATE.
  \item \textbf{Structural complexity.} Number of gates, mosaic
    complexity, reversal rate --- indicating how structurally rich the
    data is.
\end{enumerate}

\subsubsection{LLM Narrative Summary}

The JSON report is passed to an LLM (Claude API) that generates a
plain-language interpretation of the results. The LLM receives
\emph{only} the redacted summary --- never the kernel internals ---
and produces:

\begin{itemize}
  \item A one-paragraph executive summary (``Your sample shows a
    major structural transition at 93\,K with precursor activity
    beginning at 106\,K.'')
  \item Per-transition explanations in non-specialist language
  \item Flagged anomalies (``The precursor onset is unusually far
    from the transition --- this may indicate competing structural
    phases.'')
  \item Suggested follow-up measurements based on what the kernel
    found (e.g., ``The high uncertainty between 95--98\,K suggests
    a mixed-phase region; a magnetization measurement in this range
    would clarify.'')
\end{itemize}

The LLM layer is critical because most customers are not kernel
experts. They need to understand what the analysis means for their
specific work. The LLM translates structural geometry into
actionable conclusions.

\textbf{LLM security constraint:} The LLM prompt includes only the
redacted summary fields. The system prompt explicitly instructs the
LLM not to speculate about the analysis method. The LLM never sees
SEV sequences, gate structures, DSF 7-tuples, or any kernel internals.
The SES envelope is never opened for the LLM --- it receives the same
output the customer would receive, and adds interpretation on top.

\textbf{Implementation:} Claude API (Anthropic) via the existing
Anthropic SDK integration. The LLM call is a single structured prompt
with the JSON report embedded. Response is appended to the customer
report as a ``Summary'' section. Cost per analysis: approximately
\$0.01--0.05 in API tokens --- negligible relative to service pricing.

\subsection{What the Service Does NOT Expose}

The following are classified as UF-SP (UF Structural Payloads) under
SES and are never included in customer-facing output:
\begin{itemize}
  \item SEV sequences (L0 internal state)
  \item Gate boundary indices and raw gate structures (L1)
  \item Interpretive vectors, $\chi_k$, $\psi_k$ (L2)
  \item Resonance curves, URF values, gating decisions (L3)
  \item DSF 7-tuples ($D_k$, $M_k$, $U^*_k$, $B_k$, $P_k$, $C_k$,
    $R_{\text{rev},k}$) in raw form (L4)
  \item Any kernel parameters, thresholds, or algorithm details
\end{itemize}

The customer receives \emph{derived conclusions} only. The structural
geometry that produced those conclusions remains encrypted inside the
SES envelope on our infrastructure.

% ==================================================================
\section{Key Findings from Capability Proof}
% ==================================================================

\subsection{Universal Transition Detection}

The kernel detected phase transitions across:
\begin{itemize}
  \item Conventional superconductors (BCS: MgB$_2$)
  \item Unconventional superconductors (cuprate: YBCO; iron-based: FeSe/STO;
    kagome: CsV$_3$Sb$_5$)
  \item Metal--insulator transitions (Mott: VO$_2$)
  \item Ferroelectric transitions (displacive: BaTiO$_3$)
\end{itemize}

The same code path, same kernel parameters, same pipeline. No
per-material tuning.

\subsection{Precursor Detection}

In every material tested, the kernel identified structural changes
\emph{before} the nominal transition:

\begin{table}[h]
\centering
\begin{tabular}{@{}llll@{}}
\toprule
\textbf{Material} & \textbf{Known $T_c$} & \textbf{Precursor onset} &
  \textbf{Lead} \\
\midrule
YBCO & 93\,K & $\sim$106\,K & 13\,K \\
MgB$_2$ & 39\,K & $\sim$57\,K & 18\,K \\
FeSe/STO & 65\,K & $\sim$85\,K & 20\,K \\
VO$_2$ & 340\,K & $\sim$320\,K & 20\,K$^*$ \\
BaTiO$_3$ & 393\,K & $\sim$415\,K & 22\,K$^*$ \\
\bottomrule
\end{tabular}
\caption{Precursor lead distances. $^*$VO$_2$ and BaTiO$_3$ precursors
extend in the direction of the structural change (below $T_{\text{MIT}}$
and above $T_{\text{Curie}}$ respectively).}
\end{table}

This precursor detection is the service's primary differentiator. No
existing general-purpose tool provides automatic precursor onset
identification from a single transport or thermodynamic measurement.

\subsection{Universal Uncertainty Ceiling}

In every material tested, the kernel uncertainty $U_k$ peaked at
\textbf{0.667} at the transition point. This value was not programmed
--- it emerges from the structural geometry of the kernel when the
input field is maximally ambiguous. This universality across different
physical phenomena suggests the kernel is measuring a fundamental
property of structural phase transitions.

% ==================================================================
\section{Target Markets}
% ==================================================================

\subsection{Condensed Matter Physics}

\textbf{Customer:} University and national lab research groups. \\
\textbf{Use case:} Characterization of new materials, identification of
  unexpected transitions, precursor analysis. \\
\textbf{Entry point:} User facilities (PARADIM, NIST, Argonne) where
  hundreds of researchers bring samples annually.

\subsection{Pharmaceutical Polymorph Screening}

\textbf{Customer:} Pharma companies (Pfizer, Merck, Novartis, etc.). \\
\textbf{Use case:} Identifying all crystalline polymorphs of a drug
  compound. Polymorph screening is an FDA regulatory requirement.
  DSC (Differential Scanning Calorimetry) curves are the standard
  measurement --- structurally identical to $R(T)$ data. \\
\textbf{Value proposition:} Automated screening of hundreds of DSC
  curves per week. Detection of subtle polymorphic transitions that
  manual inspection misses. \\
\textbf{Revenue potential:} Highest-value segment. Regulatory
  requirement creates non-optional demand.

\subsection{Battery Materials}

\textbf{Customer:} Battery companies and research labs. \\
\textbf{Use case:} Phase transitions in electrode materials during
  charge/discharge cycling cause capacity degradation. Early detection
  of structural changes predicts cell lifetime. \\
\textbf{Data type:} Voltage vs.\ capacity curves, impedance spectra,
  cycling data.

\subsection{Semiconductor Process Monitoring}

\textbf{Customer:} Fab process engineers. \\
\textbf{Use case:} Detecting structural phase transitions in thin films
  during deposition, annealing, or doping. Quality control on novel
  materials (HfO$_2$ ferroelectrics, phase-change memory materials).

\subsection{Medical Time Series}

\textbf{Customer:} Medical device companies, hospitals. \\
\textbf{Use case:} Seizure onset detection in EEG, cardiac event
  precursors in EKG. The kernel's precursor detection capability
  translates directly to early-warning systems. \\
\textbf{Regulatory path:} FDA 510(k) or De Novo, depending on
  predicate device.

\subsection{Geological and Climate Data}

\textbf{Customer:} Oil/gas exploration, climate research. \\
\textbf{Use case:} Regime changes in seismic data, climate time series.
  Structural transitions in rock core measurements under pressure.

% ==================================================================
\section{Security Architecture}
% ==================================================================

\subsection{SES Integration}

All kernel operations are protected by the Secure Envelope System (SES)
as specified in the UF-Core Kernel Spec (see
\texttt{UF\_Core\_Kernel\_Spec\_with\_SES.pdf}). Key properties:

\begin{enumerate}
  \item \textbf{UF-SP protection.} All intermediate kernel state (SEV
    sequences, gate boundaries, interpretive vectors, resonance curves,
    DSF 7-tuples) is encrypted at rest and in transit using SCE-SIV
    authenticated encryption.
  \item \textbf{Domain parameter binding.} Each analysis is bound to a
    unique domain tuple (customer ID, data hash, timestamp, service
    version) preventing replay and cross-customer leakage.
  \item \textbf{Chain of custody.} Every analysis event is logged in an
    immutable custody chain with cryptographic hashing.
  \item \textbf{Nonce-misuse resistance.} SCE-SIV mode ensures that
    even if implementation errors cause nonce reuse, the structural
    invariants of the kernel output remain obfuscated.
  \item \textbf{Threat responses T1--T10.} The SES threat model covers
    code extraction, tampering, UF-SP exfiltration, envelope replay,
    rollback, adversarial environments, rogue tenants, domain
    manipulation, data integrity tampering, and denial of service.
\end{enumerate}

\subsection{Deployment Model}

\begin{verbatim}
Customer --> HTTPS POST (CSV) --> API Gateway --> Kernel (L0-L4)
                                                      |
                                              SES Envelope (encrypted)
                                                      |
                                              Redacted Summary
                                                      |
                                              JSON Report --> Customer
\end{verbatim}

The kernel runs exclusively on our infrastructure. No client-side
installation. No SDK. No kernel code distribution.

% ==================================================================
\section{Pricing Model}
% ==================================================================

\begin{table}[h]
\centering
\begin{tabular}{@{}llll@{}}
\toprule
\textbf{Tier} & \textbf{Customer} & \textbf{Price} & \textbf{Access} \\
\midrule
Per-analysis & Walk-in / trial & \$25--50 per curve & Web upload \\
Academic & University labs & \$3,000--5,000/yr & API, unlimited \\
Industry & Battery, semiconductor & \$25,000--50,000/yr & API + SLA \\
Enterprise & Pharma, defense & \$100,000--200,000/yr & Custom + priority \\
\bottomrule
\end{tabular}
\caption{Pricing tiers.}
\end{table}

\subsection{Pricing Rationale}

\begin{itemize}
  \item \textbf{Below beam time.} National lab beam time costs
    \$100--500/hour. A full DSF-AI analysis takes seconds.
  \item \textbf{Below simulation licenses.} Schr\"{o}dinger and similar
    computational chemistry platforms cost \$50K--200K/year per seat.
  \item \textbf{Above commodity tools.} Simple peak-finding or
    curve-fitting software is free or cheap. DSF-AI provides structural
    analysis that those tools cannot --- particularly precursor detection
    and regime classification.
  \item \textbf{Pharma premium justified.} Polymorph screening is a
    regulatory requirement with direct financial consequences (patent
    protection, FDA approval timelines). The enterprise tier reflects
    the value of automated, reliable screening at scale.
\end{itemize}

% ==================================================================
\section{Promotion Strategy}
% ==================================================================

\textbf{Constraints:} The founder is not a scientist, will not publish
papers, will not present at conferences, and has expressive dysphasia
that makes live pitching impractical. The promotion strategy must work
for a solo non-academic founder with a working product and no
connections in the target markets. Every phase below can be executed
from a laptop.

\subsection{Phase 1: The Tool Sells Itself (Cost: \$0)}

Build a public web page where anyone can upload a CSV and get a free
analysis of their first curve. No signup. No credit card. No pitch.
Just a file upload button and a result.

The page needs three things:
\begin{enumerate}
  \item The upload form
  \item The YBCO result as a worked example (``Here's what DSF-AI
    found in the most famous R(T) curve in physics --- including
    precursors that took the field a decade to discover'')
  \item A short video (screen recording, no face, no voice ---
    just the tool running on data with text annotations)
\end{enumerate}

This is the entire marketing asset. If the tool works, the result
is the pitch. If someone uploads their data and gets back something
they didn't know, they'll tell people.

\subsection{Phase 2: Seed the Right Forums (Cost: \$0)}

Post the free tool link (not a sales pitch --- the actual tool) to:
\begin{itemize}
  \item \textbf{Reddit:} r/physics, r/CondensedMatter, r/MaterialsScience,
    r/batterytech, r/pharma --- ``I built a tool that finds phase
    transitions in any measurement data. Here's what it found in YBCO.
    Try it on your data.''
  \item \textbf{Twitter/X:} Tag researchers who post R(T) curves or
    DSC data. ``Ran your published data through our structural
    analysis tool --- found precursors at [X]\,K. Full result:
    [link].'' This requires no speaking, no credentials, just results.
  \item \textbf{ResearchGate / Google Scholar alerts:} Monitor new
    papers with R(T) or DSC data. Run the kernel on their published
    figures. Email the authors: ``We ran your Figure 3 through our
    structural analysis tool and found [specific finding]. Full report
    attached. The tool is free for your first analysis: [link].''
  \item \textbf{Hacker News:} ``Show HN: Tool that finds phase
    transitions in any time series data --- no physics knowledge
    required.'' HN loves domain-agnostic tools. If it hits the front
    page, thousands of technical users see it in a day.
\end{itemize}

None of this requires being a scientist, giving a talk, or having
credentials. It requires having a tool that works and showing results.

\subsection{Phase 3: YouTube / Content Without Speaking}

Create short (2--3 minute) screen-recording videos:
\begin{itemize}
  \item ``Uploading YBCO data --- watch DSF-AI find the pseudogap''
  \item ``Can a non-physics tool find a ferroelectric transition?''
  \item ``Battery cycling data: where does the electrode start
    degrading?''
\end{itemize}

No voiceover needed. Text overlays, screen capture, background music.
These are technical demos, not vlogs. The audience is researchers
and engineers who search for tools. YouTube SEO on terms like
``phase transition detection,'' ``DSC analysis tool,''
``R(T) analysis software'' puts the videos in front of people
actively looking for solutions.

\subsection{Phase 4: Cold Outreach via Results (Cost: \$0)}

This is the highest-conversion strategy and requires no credentials:
\begin{enumerate}
  \item Find a published paper with measurement data (R(T), DSC,
    impedance, etc.)
  \item Digitize the data from the figure
  \item Run the kernel
  \item If the kernel finds something the paper didn't mention ---
    a precursor, a subtle transition, a regime boundary --- email
    the corresponding author with the result attached
  \item The email is not a sales pitch. It is: ``I ran your Figure 3
    data through a structural analysis tool. It found [X] at [Y].
    Thought you'd want to know. The tool is at [link] if you want
    to try your other datasets.''
\end{enumerate}

This works because scientists care about results, not credentials.
If the kernel finds something real in their data, they will use
the tool. If it finds nothing interesting, nothing is lost.

\textbf{Volume target:} 10 papers per week. At even a 5\% conversion
rate, that's 2 new users per month from cold outreach alone, each of
whom tells their lab group.

\subsection{Phase 5: Institutional Licensing (When Revenue Justifies)}

Once individual researchers are using the tool and generating word of
mouth, approach institutions for site licenses:
\begin{itemize}
  \item University materials science departments
  \item National lab user facilities
  \item Battery company R\&D groups
  \item Pharma analytical chemistry groups
\end{itemize}

This is the only phase that may require a business partner or sales
representative. By this point, the tool has users and results to
point to. The partner's job is negotiating contracts, not explaining
the product.

\subsection{What NOT to Do}

\begin{itemize}
  \item Do not apply to accelerators, incubators, or pitch competitions.
    They optimize for presentation skills, not product quality.
  \item Do not hire a marketing agency. They will produce generic
    content that means nothing to the target audience.
  \item Do not attend conferences until invited. Invitations follow
    results, not the other way around.
  \item Do not explain the method. Ever. Show the result. The method
    is a trade secret and explaining it adds nothing to the sale.
\end{itemize}

% ==================================================================
\section{Intellectual Property}
% ==================================================================

\begin{itemize}
  \item The UF-Core kernel (L0--L4) is a \textbf{trade secret}.
    It is never distributed, open-sourced, or exposed through the API.
  \item The SES/SCE encryption layer is specified in published documents
    but the implementation details of the structural engine keystream
    generation remain proprietary.
  \item Customer data is encrypted in transit and at rest. Customer data
    is not used for any purpose other than the requested analysis.
    Customer data is deleted after analysis unless the customer requests
    retention.
  \item The DSF 7-tuple field definitions and their structural
    interpretation are trade secrets.
  \item Published results (transition temperatures, precursor onsets,
    regime maps) are factual observations and do not expose trade secrets.
\end{itemize}

% ==================================================================
\section{Technical Implementation}
% ==================================================================

\subsection{Architecture}

\begin{verbatim}
[Browser]                    [Backend (AWS ECS)]
    |                              |
    |-- POST /analyze (CSV) ------>|
    |                              |-- Parse CSV
    |                              |-- build_field_series()
    |                              |-- run_kernel() [L0-L4]
    |                              |-- SES envelope (encrypt internals)
    |                              |-- Extract redacted summary
    |                              |-- Claude API call (narrative)
    |                              |-- Assemble report
    |<---- JSON + narrative -------|
\end{verbatim}

\subsection{Backend Components}

\begin{enumerate}
  \item \textbf{CSV Parser.} Accepts two-column CSV. Validates
    numeric data, strips headers, handles common formats (comma,
    tab, semicolon delimited). Rejects files $>$ 10\,MB or
    $>$ 100{,}000 rows (configurable per tier).

  \item \textbf{Kernel Runner.} Calls \texttt{build\_field\_series()}
    and \texttt{run\_kernel()} from
    \texttt{tools/derive\_sppu\_weights.py}. Returns DSF list,
    gate boundaries, and gate count. This is the trade secret
    component --- runs in an isolated container with no external
    network access.

  \item \textbf{SES Wrapper.} Encrypts all kernel internals (SEV
    sequences, gate structures, DSF 7-tuples) into an SES envelope
    using SCE-SIV. The envelope is stored server-side for audit
    purposes. Only the redacted summary leaves the secure boundary.

  \item \textbf{Report Generator.} Extracts from the kernel output:
    \begin{itemize}
      \item Transition temperatures (from gate boundary indices
        mapped back to stimulus values)
      \item Regime labels (from L2 interpretations)
      \item Precursor onset (first gate where L0 $\sigma$ exceeds
        baseline by $>$ 10$\times$)
      \item Uncertainty peak (maximum $U_k$ and its location)
      \item Structural complexity (gate count, reversal rate)
    \end{itemize}
    These are factual summaries that expose no kernel internals.

  \item \textbf{LLM Narrator.} Sends the redacted summary to
    Claude API with a structured system prompt:
    \begin{itemize}
      \item ``You are a materials analysis assistant. You receive
        structural analysis results from a proprietary tool. Explain
        the findings in clear, non-specialist language. Do not
        speculate about the analysis method. Focus on what was found
        and what it means for the user's material.''
      \item The user's optional context (``this is a battery cathode
        material'' or ``this is a pharmaceutical compound'') is
        included if provided, allowing domain-appropriate language.
    \end{itemize}

  \item \textbf{Authentication and metering.} API key per customer.
    Rate limiting per tier. Usage tracking for billing. Standard
    implementation (AWS API Gateway + DynamoDB counters).
\end{enumerate}

\subsection{Infrastructure}

\begin{itemize}
  \item \textbf{Compute:} AWS ECS Fargate (same infrastructure as
    the existing TFE deployment). Kernel runs in an isolated task
    with no outbound network access.
  \item \textbf{Storage:} SES envelopes stored in S3 with
    server-side encryption (SSE-KMS). Customer CSVs deleted after
    analysis unless retention is requested.
  \item \textbf{Cost per analysis:} Approximately \$0.001 compute
    + \$0.01--0.05 LLM API = under \$0.10 total cost per curve.
    At \$25/curve walk-in pricing, gross margin exceeds 99\%.
\end{itemize}

\subsection{Web Frontend}

A single-page application with:
\begin{enumerate}
  \item File upload dropzone
  \item Optional context field (``What kind of data is this?'')
  \item Results display: regime map visualization, transition table,
    narrative summary
  \item Download buttons for JSON and PDF report
  \item ``First analysis free'' for unauthenticated users; API key
    required after that
\end{enumerate}

No framework dependencies beyond what is needed. The frontend is a
sales tool --- it must load fast, work on mobile, and produce a
shareable result URL.

\subsection{Development Roadmap}

\begin{enumerate}
  \item \textbf{Phase 1: Working demo.} CSV upload $\to$ kernel
    $\to$ JSON report $\to$ LLM narrative. No authentication, no
    billing. Deployed on existing ECS infrastructure. This is the
    ``free first analysis'' tool used for promotion.

  \item \textbf{Phase 2: Authentication and billing.} API keys,
    Stripe integration, tiered rate limiting. This converts free
    users to paying customers.

  \item \textbf{Phase 3: Batch processing.} ZIP upload for
    high-throughput screening. Essential for pharma polymorph
    and battery cycling use cases.

  \item \textbf{Phase 4: Multi-column and comparison.} Accept
    datasets with multiple measurement columns or multiple files
    for the same sample under different conditions. Cross-correlate
    structural findings.

  \item \textbf{Phase 5: Historical database.} Store (encrypted)
    analysis results to enable ``has this structural signature been
    seen before?'' queries across a customer's dataset history.
    This is the lock-in feature --- the more data a customer runs
    through the system, the more valuable their history becomes.
\end{enumerate}

% ==================================================================
\section{Appendix: Detailed Results}
% ==================================================================

\subsection{YBCO --- YBa$_2$Cu$_3$O$_{7-\delta}$}

Data source: Wu et al., PRL 58, 908 (1987), Figure 3.

\begin{itemize}
  \item \textbf{Kernel segmentation:} 31 gates.
    Gate 0 (4--77\,K): TRANSITIONAL. Gates 1--29 (78--106\,K): VOLATILE.
    Gate 30 (107--300\,K): DEGENERATE.
  \item \textbf{Transition:} $D_k$ reversal at 92\,K (known $T_c = 93$\,K).
  \item \textbf{Precursors:} L0 $\sigma$ begins rising at $\sim$105\,K.
    VOLATILE regime extends to 106\,K. $U_k$ peaks at 0.660 at 98\,K.
  \item \textbf{Interpretation:} The kernel detects the pseudogap /
    superconducting fluctuation regime extending $\sim$13\,K above $T_c$,
    consistent with findings that required a decade of specialized
    measurements to establish.
\end{itemize}

\subsection{MgB$_2$}

Data source: Nagamatsu et al., Nature 410, 63 (2001), Figure 2.

\begin{itemize}
  \item \textbf{Kernel segmentation:} 23 gates.
    STABLE (2--36\,K). VOLATILE (36--57\,K). DEGENERATE (57--300\,K).
  \item \textbf{Transition:} $D_k$ reversals at 37--39\,K (known $T_c = 39$\,K).
  \item \textbf{Precursors:} VOLATILE regime extends to 57\,K, 18\,K above $T_c$.
  \item \textbf{Note:} MgB$_2$ is a conventional BCS superconductor.
    Precursors above $T_c$ in a BCS system are unexpected and may
    indicate superconducting fluctuation effects previously
    underappreciated in this material.
\end{itemize}

\subsection{CsV$_3$Sb$_5$}

Data source: Ortiz et al., PRL 125, 247002 (2020), Figure 3.

\begin{itemize}
  \item \textbf{Kernel segmentation:} 26 gates.
    STABLE (0--2\,K). VOLATILE (2--26\,K). DEGENERATE (26--300\,K).
  \item \textbf{SC transition:} $D_k$ reversals at 4--5\,K (known $T_{c,\text{SC}} = 2.5$\,K).
  \item \textbf{CDW transition:} L0 $\sigma$ shows elevated values
    near 90--98\,K (known $T_{\text{CDW}} = 94$\,K) but within the
    DEGENERATE regime. The CDW manifests as a slope change in $R(T)$,
    not a resistance drop, making it a subtler structural event.
  \item \textbf{Significance:} This is a kagome superconductor ---
    the same lattice geometry as the Cu$_3$O$_2$ target material in
    the RTSC project.
\end{itemize}

\subsection{VO$_2$}

Data source: Morin, PRL 3, 34 (1959); composite modern data.

\begin{itemize}
  \item \textbf{Kernel segmentation:} 44 gates (most of any material tested).
    DEGENERATE (200--320\,K). VOLATILE (320--362\,K). TRANSITIONAL (362--400\,K).
  \item \textbf{Transition:} $D_k$ reversals at 337--342\,K
    (known $T_{\text{MIT}} = 340$\,K). $U_k$ peaks at 0.667 at 341\,K.
  \item \textbf{Precursors:} VOLATILE regime begins at 320\,K,
    20\,K below $T_{\text{MIT}}$.
  \item \textbf{Significance:} This is NOT a superconductor. It is a
    Mott metal--insulator transition involving a 4--5 order of magnitude
    resistivity change. The kernel detected it using the identical code
    path as the superconductor analyses, demonstrating true
    domain-agnosticism.
\end{itemize}

\subsection{BaTiO$_3$}

Data source: Merz, Physical Review 76, 1221 (1949); composite data.

\begin{itemize}
  \item \textbf{Measurement type:} Dielectric constant vs.\ temperature
    (NOT resistance).
  \item \textbf{Kernel segmentation:} 22 gates.
    DEGENERATE (300--395\,K). VOLATILE (395--415\,K). DEGENERATE (415--500\,K).
  \item \textbf{Transition:} $D_k$ reversal at 396\,K
    (known $T_{\text{Curie}} = 393$\,K).
  \item \textbf{Significance:} The kernel analyzed a completely
    different physical quantity (permittivity, not resistance) and
    found the ferroelectric Curie point. This confirms that the kernel
    operates on structural field geometry, not on any property-specific
    feature.
\end{itemize}

\subsection{FeSe/STO}

Data source: Wang et al., Chinese Physics Letters 29, 037402 (2012).

\begin{itemize}
  \item \textbf{Kernel segmentation:} 27 gates.
    TRANSITIONAL (4--60\,K). VOLATILE (60--85\,K). DEGENERATE (85--300\,K).
  \item \textbf{Transition:} $D_k$ reversals at 61--63\,K (known $T_c \approx 65$\,K).
    $U_k$ peaks at 0.667 at 62\,K.
  \item \textbf{Precursors:} VOLATILE regime extends to 85\,K,
    20\,K above $T_c$.
  \item \textbf{Significance:} Iron-based superconductors have a
    different pairing mechanism than cuprates. The kernel finds the
    same structural signature pattern regardless of the underlying
    physics.
\end{itemize}

% ==================================================================
\section{Glossary}
% ==================================================================

Terms used throughout this document, ordered by dependency (foundational
terms first).

\begin{description}

  \item[UF-Core] \textbf{Universal Field Core.} The five-layer
    deterministic analysis engine (L0--L4) that processes any
    measurement-vs-stimulus series into structural geometry. Domain-agnostic
    --- it does not know what the data represents. Trade secret.

  \item[SEV (Structural Event Value)] The L0 output. For each data
    point, L0 computes four quantities: $F_{\text{norm}}$ (normalized
    field value), $dF$ (first derivative / rate of change), $\sigma$
    (structural change rate --- how fast the local geometry is
    shifting), and $\kappa$ (curvature). These are the raw structural
    features the rest of the pipeline builds on.

  \item[Gate] A contiguous segment of the data series where the
    structural character is internally consistent. L1 finds the
    boundaries between gates --- the points where the data's
    structural behavior changes. A phase transition always appears
    as a gate boundary. The number of gates indicates how structurally
    complex the data is.

  \item[Regime] L2's classification of each gate's structural
    character. Four types:
    \begin{description}
      \item[STABLE] Low $\sigma$, low $U_k$. The data is structurally
        quiet --- nothing is changing. Example: zero-resistance state
        in a superconductor below $T_c$.
      \item[VOLATILE] High $\sigma$, variable $U_k$. The data is
        structurally active --- something is happening or about to
        happen. Phase transitions always occur within VOLATILE regions.
      \item[TRANSITIONAL] Intermediate character. The data is moving
        between structural states but hasn't committed.
      \item[DEGENERATE] Structurally featureless. The data has a
        consistent trend with no internal structure worth segmenting.
        Example: linear metallic resistance above all transitions.
    \end{description}

  \item[DSF (Dynamic Structural Field)] The L4 output --- a
    7-tuple that captures the complete structural geometry of each gate:
    \begin{description}
      \item[$D_k$ (Direction)] Which way the structural field points
        ($+1$, $0$, or $-1$). A reversal of $D_k$ between adjacent
        gates marks a structural direction change --- the kernel's
        way of saying ``something fundamentally changed here.''
      \item[$M_k$ (Momentum)] Rate of structural change. How fast the
        field geometry is evolving.
      \item[$U^*_k$ (Uncertainty)] How confident the kernel is in its
        structural assessment. Peaks at $\sim$0.667 at phase transitions
        --- maximum structural ambiguity.
      \item[$B_k$ (Breathing)] Structural freedom / expansion. How
        much room the field has to move.
      \item[$P_k$ (Pressure)] Transition sensitivity. Spikes when the
        structural state is under stress / about to change.
      \item[$C_k$ (Complexity)] Mosaic divergence. How many distinct
        structural projections exist within a gate.
      \item[$R_{\text{rev},k}$ (Reversal)] How often the structural
        direction flips within a gate. High reversals = unstable signal.
    \end{description}

  \item[Precursor] A structural change detected by the kernel
    \emph{before} the nominal transition. Manifests as rising $\sigma$
    in L0 or the onset of a VOLATILE regime above/before the known
    transition temperature. The kernel's primary novel capability.

  \item[SES (Secure Envelope System)] The encryption and access
    control layer that protects kernel internals. All intermediate
    state (SEV sequences, gate structures, DSF 7-tuples) is encrypted
    using SCE-SIV and stored in an SES envelope. Only redacted
    summaries (transition temperatures, regime labels, precursor
    onsets) leave the secure boundary.

  \item[SCE (Structural Cryptographic Engine)] The cryptographic
    backend used by SES. Provides deterministic authenticated
    encryption (SCE-SIV mode) with nonce-misuse resistance.
    Implementation details are proprietary.

  \item[UF-SP (UF Structural Payload)] The classification for all
    kernel internals that must never be exposed to customers.
    Includes SEV sequences, gate structures, interpretive vectors,
    resonance curves, and raw DSF 7-tuples.

  \item[LLM Narrator] The Claude API integration that translates the
    kernel's structured JSON output into plain-language narrative.
    Receives only redacted summaries --- never kernel internals.

  \item[CSV] Comma-Separated Values. The input format. Two columns:
    stimulus (e.g., temperature) and measurement (e.g., resistance).

  \item[Phase transition] Any point where a material's physical
    properties change abruptly due to a change in its internal
    structure. Examples: superconducting transition (resistance
    drops to zero), metal--insulator transition (resistance jumps
    by orders of magnitude), ferroelectric transition (dielectric
    constant peaks). The kernel detects all of these as structural
    geometry changes without knowing which physical phenomenon is
    involved.

  \item[$T_c$] Critical temperature. The temperature at which a
    phase transition occurs. Different subscripts for different
    transitions: $T_{c,\text{SC}}$ for superconducting,
    $T_{\text{MIT}}$ for metal--insulator, $T_{\text{Curie}}$ for
    ferroelectric, $T_{\text{CDW}}$ for charge density wave.

  \item[R(T)] Resistance as a function of temperature. The most
    common measurement type for transport characterization.

  \item[DSC (Differential Scanning Calorimetry)] A measurement of
    heat flow vs.\ temperature. Standard technique for polymorph
    screening in pharmaceuticals. Structurally identical to R(T)
    from the kernel's perspective --- it's just another
    measurement-vs-stimulus curve.

\end{description}

% ==================================================================
\section{Hosting and Domain Architecture}
% ==================================================================

\subsection{Domain Strategy}

The DSF-AI service runs on its own domain, completely separate from
the TFE trading infrastructure. The two systems share no public-facing
resources.

\begin{itemize}
  \item \textbf{Service domain:} e.g., \texttt{dsf-ai.com} or
    \texttt{structuralfield.ai} (to be registered).
  \item \textbf{API endpoint:} \texttt{api.dsf-ai.com/v1/analyze}
  \item \textbf{Web app:} \texttt{dsf-ai.com} (the upload page /
    marketing / results)
  \item \textbf{TFE domain:} Unchanged. No cross-linking, no shared
    authentication, no shared infrastructure visible to customers.
\end{itemize}

\subsection{Why AWS (Not Azure)}

\begin{itemize}
  \item TFE already runs on AWS ECS Fargate. The kernel code, Python
    environment, and SES libraries are already containerized and
    deployed there. Using AWS means zero re-platforming.
  \item AWS Secrets Manager already holds the encryption keys (SES
    root key provider). Moving to Azure would require migrating the
    key management infrastructure.
  \item Claude API (Anthropic) works from anywhere --- no cloud
    vendor dependency for the LLM layer.
  \item AWS free tier covers the initial deployment cost for low
    traffic. ECS Fargate pricing is per-second --- idle costs are
    near zero.
\end{itemize}

\subsection{AWS Architecture}

\begin{verbatim}
[Route 53]  dsf-ai.com
     |
[CloudFront]  CDN for static assets + SSL termination
     |
     +-- /           --> [S3]  Static website (HTML/CSS/JS)
     |
     +-- /api/v1/*   --> [API Gateway]
                              |
                         [Lambda or ECS Fargate]
                              |
                         +-- CSV Parser
                         +-- Kernel Runner (isolated, no network)
                         +-- SES Wrapper
                         +-- Report Generator
                         +-- Claude API call (outbound only)
                              |
                         [S3]  SES envelope storage (SSE-KMS)
                         [DynamoDB]  API keys, usage metering
\end{verbatim}

\textbf{Two deployment options for the backend:}

\begin{enumerate}
  \item \textbf{Lambda (simpler, cheaper at low volume).} Each
    analysis request is a single Lambda invocation. The kernel runs
    in under 5 seconds for typical datasets. Lambda's 15-minute
    timeout is more than sufficient. Cold start adds $\sim$2 seconds
    for the Python runtime. Cost: \$0.0002 per invocation + compute
    time. Scales to zero when idle.

  \item \textbf{ECS Fargate (better for sustained traffic).} Same
    container image as the existing TFE deployment. Always-on task
    with auto-scaling. Better for batch processing (Phase 3 of
    roadmap). Cost: $\sim$\$30/month for a single always-on task.
\end{enumerate}

\textbf{Recommendation:} Start with Lambda for Phase 1 (free demo
tool). Switch to ECS Fargate when batch processing or sustained
traffic justifies it.

\subsection{SSL and Security}

\begin{itemize}
  \item SSL certificate via AWS Certificate Manager (free for
    CloudFront and API Gateway).
  \item All traffic HTTPS. No HTTP fallback.
  \item API Gateway handles rate limiting, API key validation,
    and request throttling.
  \item Kernel container has no outbound network access except
    the Claude API endpoint (whitelisted).
  \item Customer CSV data is encrypted in transit (TLS 1.3) and
    at rest (S3 SSE-KMS). Deleted after analysis unless customer
    opts in to retention.
\end{itemize}

\subsection{Cost Estimate (Monthly)}

\begin{table}[h]
\centering
\begin{tabular}{@{}llll@{}}
\toprule
\textbf{Component} & \textbf{Low traffic} & \textbf{Medium} &
  \textbf{High} \\
 & (50 analyses/mo) & (500/mo) & (5{,}000/mo) \\
\midrule
Route 53 & \$0.50 & \$0.50 & \$0.50 \\
CloudFront & \$0 (free tier) & \$1 & \$5 \\
S3 (static site) & \$0.03 & \$0.03 & \$0.03 \\
Lambda compute & \$0.01 & \$0.10 & \$1.00 \\
API Gateway & \$0.18 & \$1.75 & \$17.50 \\
S3 (envelopes) & \$0.01 & \$0.12 & \$1.15 \\
DynamoDB & \$0 (free tier) & \$0 & \$1.25 \\
Claude API (LLM) & \$2.50 & \$25 & \$250 \\
\midrule
\textbf{Total} & \textbf{\$3.23} & \textbf{\$28.50} & \textbf{\$276.43} \\
\bottomrule
\end{tabular}
\caption{Monthly hosting cost estimates. Claude API dominates at scale.
At \$25/analysis, 50 analyses/mo = \$1{,}250 revenue vs.\ \$3.23 cost.}
\end{table}

% ==================================================================
\section{Website Creation}
% ==================================================================

\subsection{The Problem with Previous Approaches}

Building custom web applications from scratch (as with the TFE
Next.js dashboard) is painful, slow, and fragile. The DSF-AI
service page does not need a custom web framework. It needs:

\begin{enumerate}
  \item A file upload button
  \item A results display
  \item A payment form (later)
\end{enumerate}

That's it. No dashboard. No real-time updates. No user accounts
(initially). No database-backed UI.

\subsection{Recommended Approach: Static Site + API}

The website is a \textbf{static HTML/CSS/JS page} hosted on S3 +
CloudFront. No server-side rendering. No Next.js. No React (unless
desired for the results display). The page calls the API endpoint
and displays the result.

\textbf{Why static:}
\begin{itemize}
  \item Zero server maintenance. S3 hosting costs pennies.
  \item Loads instantly (no SSR cold start, no Node.js runtime).
  \item Cannot break. There is no server-side code to crash.
  \item Deploys in seconds (\texttt{aws s3 sync}).
  \item Works on every device and browser without framework quirks.
\end{itemize}

\textbf{What the page looks like:}
\begin{enumerate}
  \item \textbf{Hero section:} ``Upload your measurement data.
    Get structural analysis in 60 seconds.'' One sentence. One
    upload button.
  \item \textbf{Worked example:} The YBCO result, pre-rendered.
    Shows exactly what a customer gets. ``Here's what DSF-AI found
    in the most famous R(T) curve in physics.''
  \item \textbf{Upload form:} Drag-and-drop CSV. Optional text
    field: ``What kind of data is this?'' Submit button.
  \item \textbf{Results area:} Appears after analysis completes.
    Shows regime map, transition table, narrative summary.
    Download buttons for JSON and PDF.
  \item \textbf{Footer:} Pricing tiers. ``First analysis free.''
    Link to API documentation (a single Markdown page).
\end{enumerate}

\subsection{Build Options (Ranked by Pain Level)}

\begin{enumerate}
  \item \textbf{Claude builds it (lowest pain).} Claude Code
    generates the entire static site --- HTML, CSS, JS --- as files
    in the repo. No framework, no build tools, no npm. The JS
    makes a \texttt{fetch()} call to the API endpoint and renders
    the result. Claude can iterate on the design in conversation.
    Deployment: \texttt{aws s3 sync ./site s3://dsf-ai-site}.

    \textit{This is the recommended approach for Phase 1.}

  \item \textbf{Astro or 11ty (low pain).} Lightweight static site
    generators. Write content in Markdown, output static HTML.
    Minimal tooling. Good if the site grows beyond a single page
    (e.g., documentation, blog posts about results).

  \item \textbf{Vercel / Netlify (medium pain).} One-click deploy
    from a Git repo. Handles SSL, CDN, and preview deploys
    automatically. Adds a vendor dependency but removes all
    infrastructure management. Free tier is generous.

    \textit{Fallback option if AWS static hosting feels like too
    much configuration.}

  \item \textbf{Next.js / React (high pain --- avoid).} The TFE
    web dashboard uses Next.js 16 and it has been a source of
    ongoing friction. For a single-page upload tool, a full React
    framework is massive overkill. Do not use unless the site
    eventually needs authenticated user dashboards with complex
    state management (Phase 5 of the roadmap at earliest).
\end{enumerate}

% ==================================================================
\section{Nanoparticle Property Screener}
% ==================================================================

\subsection{Overview}

In addition to the CSV-upload structural analysis, the service offers
a \textbf{nanoparticle property calculator} that predicts cluster
properties from geometry alone. No DFT. No simulation. Milliseconds
per prediction.

This capability was validated on May 5, 2026 against 63 published
experimental measurements across 7 elements (Au, Ag, Cu, Fe, Co, Ni,
Mn), 4 property types, cluster sizes $N = 3$ to $N = 700$ atoms, and
5 distinct geometries. Average error: 6.5\%. Zero free parameters.
Zero exclusions.

\subsection{Properties Predicted}

For any d-metal cluster of size $N$:

\begin{enumerate}
  \item \textbf{Magnetic moment} ($\mu_B$/atom). Validated for Fe, Co,
    Ni (all sizes 13--700), Mn (5--19). Fe$_{13}$ and Co$_{13}$ are
    \textbf{exact} (0.0\% error).

  \item \textbf{Electron affinity} (eV). Validated for Au ($N = 3$--20),
    Ag ($N = 3$--10), Cu ($N = 5$--10). Average error 9\%.

  \item \textbf{Seebeck coefficient} ($\mu$V/K). Validated for
    Au$_{13}$ cubic superlattice (87.6 vs.\ 93 measured, 6\% off).
    Predictions available for all 23 d-metals at all cluster sizes
    and both cubic/hexagonal lattice symmetries (240 predictions).

  \item \textbf{HOMO-LUMO gap} (eV). Validated for Au$_{13}$
    (0.648 vs.\ 0.650 measured, 0.2\% off). Predictions for all
    d-metals.
\end{enumerate}

\subsection{API Endpoint}

\begin{verbatim}
POST /api/v1/cluster
{
  "element": "Au",
  "N_atoms": 13,
  "temperature_K": 300,       // optional, default 300
  "lattice": "cubic"          // optional, for Seebeck
}

Response:
{
  "element": "Au",
  "N_atoms": 13,
  "magnetic_moment_uB": 0.0,
  "electron_affinity_eV": 3.71,
  "seebeck_uV_K": 87.6,
  "homo_lumo_gap_eV": 0.648,
  "geometry": "icosahedral",
  "d_band_screening": 0.392,
  "disruption_active": false,
  "validation_note": "Validated against 63 experiments, 6.5% avg error"
}
\end{verbatim}

\subsection{Batch Screening Endpoint}

\begin{verbatim}
POST /api/v1/cluster/screen
{
  "elements": ["Au", "Ag", "Cu", "Pt", "Pd", "Ni", "Fe", "Co"],
  "N_atoms": [13, 55, 147],
  "constraints": {
    "moment_min_uB": 1.0,
    "seebeck_min_uV_K": 50,
    "EA_min_eV": 3.0
  }
}

Response:
{
  "candidates": [
    {"element": "Co", "N": 13, "moment": 2.0, "S": 53.1, "EA": 3.20},
    {"element": "Ir", "N": 13, "moment": 2.0, "S": 52.6, "EA": 3.61},
    ...
  ],
  "screened": 24,
  "matched": 4,
  "compute_time_ms": 12
}
\end{verbatim}

This endpoint replaces months of DFT screening. A catalyst designer
specifies target properties and gets a ranked shortlist in milliseconds.
The shortlist then goes to DFT verification (top 5--10 candidates
instead of hundreds).

\subsection{Thermocouple Pair Finder}

\begin{verbatim}
POST /api/v1/cluster/thermocouple
{
  "N_atoms": 13,
  "min_delta_S": 80
}

Response:
{
  "pairs": [
    {"A": "Fe13", "B": "Co13", "delta_S": 97, "S_A": 47.6, "S_B": -50.9},
    ...
  ],
  "note": "Fe13/Co13 doubles a Type K thermocouple from cheap metals"
}
\end{verbatim}

\subsection{d-Band Design Rule}

The screener includes a universal design rule discovered during
validation: the d-band screening factor
$|{\epsilon_d}|/(|{\epsilon_d}| + E_f)$ determines whether a
cluster's superatom shell structure is clean or disrupted.
The critical threshold equals $\sin^2\theta$ where
$\theta = \arctan(1/\varphi)$ and $\varphi$ is the golden ratio.

\begin{itemize}
  \item \textbf{Above threshold:} d-electrons are spectators.
    Superatom shell model works cleanly. Au, Ag, Pd, Pt, Rh.
  \item \textbf{At/below threshold:} d-electrons disrupt the
    superatom shell. Electron affinity is reduced. Cu (exactly at
    threshold), Sc, Ti, V.
\end{itemize}

This is reported in the API response as \texttt{disruption\_active}
and allows customers to immediately know whether standard superatom
models apply to their material.

\subsection{What the Screener Does NOT Expose}

The geometric engine behind the screener is a trade secret.
Customers receive \textbf{property values only} --- never the
formulas, angles, projections, or coupling constants that produce
them. The response looks like any computational chemistry result.
The fact that it takes milliseconds instead of hours is the only
visible clue that the method is different.

\subsection{Implementation}

The screener is a single Python module
(\texttt{tools/ufcp\_nanoparticle\_screener.py} and supporting files)
that runs entirely in memory. No external dependencies beyond NumPy.
No database. No iterative solver. Each prediction is a direct
algebraic evaluation --- approximately 10 floating-point operations
per property per cluster.

\begin{itemize}
  \item \textbf{Compute cost:} Effectively zero. A single Lambda
    invocation can evaluate 10{,}000 clusters in under 1 second.
  \item \textbf{No GPU required.} No matrix operations. No
    optimization loops.
  \item \textbf{Deployment:} Same Lambda/ECS infrastructure as the
    CSV analysis endpoint. Separate route, same container.
\end{itemize}

% ==================================================================
\section{Weekend Implementation Plan}
% ==================================================================

Build and deploy the complete service in one weekend session.
Everything below uses existing infrastructure (AWS account, ECS
setup, Python environment) and existing code (kernel, screener).

\subsection{Saturday: Backend + API}

\begin{enumerate}
  \item \textbf{Create \texttt{dsf\_ai\_service/} directory} in repo.
    Contains: \texttt{app.py} (Flask/FastAPI), \texttt{kernel\_runner.py},
    \texttt{cluster\_screener.py}, \texttt{report\_generator.py},
    \texttt{narrator.py} (Claude API wrapper).

  \item \textbf{Wire up three endpoints:}
    \begin{itemize}
      \item \texttt{POST /api/v1/analyze} --- CSV upload $\to$ kernel
        $\to$ JSON report $\to$ Claude narrative
      \item \texttt{POST /api/v1/cluster} --- element + N $\to$
        screener $\to$ properties JSON
      \item \texttt{POST /api/v1/cluster/screen} --- batch screening
        with constraints
    \end{itemize}

  \item \textbf{Containerize.} Dockerfile based on existing TFE image.
    Add Flask/FastAPI, strip TFE-specific code. Kernel and screener
    modules only.

  \item \textbf{Deploy to ECS Fargate} on existing AWS account.
    Single task, 0.5 vCPU, 1 GB RAM (kernel is lightweight). No
    public endpoint yet --- test via \texttt{curl} from Codespace.

  \item \textbf{Test:} Upload YBCO CSV, verify kernel output matches
    previous results. Call cluster endpoint for Au$_{13}$, verify
    87.6 $\mu$V/K.
\end{enumerate}

\subsection{Sunday: Frontend + Go Live}

\begin{enumerate}
  \item \textbf{Static site} (HTML/CSS/JS, no framework):
    \begin{itemize}
      \item Hero: ``Upload measurement data. Get structural analysis
        in 60 seconds.''
      \item Tab 1: CSV upload (phase transition analysis)
      \item Tab 2: Cluster screener (element picker + size slider)
      \item Worked examples: YBCO R(T) result, Au$_{13}$ cluster
        properties
      \item Results display: regime map, transition table, narrative,
        cluster property cards
    \end{itemize}

  \item \textbf{Upload to S3}, front with CloudFront. SSL via ACM.

  \item \textbf{API Gateway} in front of ECS. Rate limiting: 10
    requests/minute for unauthenticated, unlimited for API key
    holders. First analysis free (no auth required).

  \item \textbf{Connect domain} (register Saturday if not done).

  \item \textbf{Test end-to-end:} Upload CSV from browser, see
    result. Pick Au/13 from cluster screener, see properties.
    Share URL with test user.
\end{enumerate}

\subsection{What Ships}

By Sunday evening:
\begin{itemize}
  \item Live website at \texttt{dsf-ai.com} (or chosen domain)
  \item CSV structural analysis (kernel, 6 validated materials)
  \item Nanoparticle screener (23 elements, 63 validated predictions)
  \item Claude narrative for CSV results
  \item First analysis free, no signup
  \item Shareable result URLs
\end{itemize}

\subsection{What Does NOT Ship Yet}

\begin{itemize}
  \item Authentication / API keys (Phase 2)
  \item Stripe billing (Phase 2)
  \item Batch ZIP upload (Phase 3)
  \item SES envelope storage (Phase 2 --- kernel internals are
    protected by not being in the response, but formal encryption
    comes later)
  \item PDF report download (nice-to-have, not launch-critical)
\end{itemize}

\subsection{Domain Registration}


\begin{itemize}
  \item Register domain via AWS Route 53 (\$12/year for .com) or
    any registrar (Namecheap, Cloudflare Registrar).
  \item Suggested names (check availability at registration time):
    \begin{itemize}
      \item \texttt{dsf-ai.com}
      \item \texttt{structuralfield.ai} (.ai domains $\sim$\$80/year)
      \item \texttt{phasedetect.com}
      \item \texttt{structuralanalysis.io}
    \end{itemize}
  \item Point DNS to CloudFront distribution. SSL certificate via
    ACM (free, auto-renewing).
\end{itemize}

\subsection{Timeline}

\begin{table}[h]
\centering
\begin{tabular}{@{}lll@{}}
\toprule
\textbf{Task} & \textbf{Effort} & \textbf{Who} \\
\midrule
Register domain & 10 minutes & Joseph \\
Static site (HTML/CSS/JS) & 1 session & Claude \\
API endpoint (Lambda) & 1 session & Claude \\
S3 + CloudFront setup & 30 minutes & Claude (via CLI) \\
Connect domain + SSL & 15 minutes & Claude (via CLI) \\
YBCO worked example & Already done & --- \\
\midrule
\textbf{Total to live site} & \textbf{1--2 sessions} & \\
\bottomrule
\end{tabular}
\caption{From zero to live website. The kernel and worked example
already exist. The remaining work is plumbing.}
\end{table}

% ==================================================================
\section{Deployment Status (May 11, 2026)}
% ==================================================================

The service is \textbf{live} at \texttt{https://dsf-ai.com}.

\subsection{What Shipped}

\begin{enumerate}
  \item \textbf{CSV Structural Analysis.} Upload any two-column CSV.
    Returns transitions, precursors, regime map, uncertainty profile.
    Tested on 7 materials across 4 physical phenomena plus
    pharmaceutical DSC data (5/5 transitions detected).

  \item \textbf{Nanoparticle Cluster Screener.} Predicts magnetic moment,
    electron affinity, Seebeck coefficient, HOMO-LUMO gap for any
    d-metal cluster. 23 elements, 5 cluster sizes. Always free.

  \item \textbf{Thermocouple Pair Finder.} Ranks all element pairs by
    Seebeck difference. Auto-updates on filter change.

  \item \textbf{Validation Evidence Page.} 63 experiments, 4 property
    types, 6 phase transitions, all sources cited. Zero free parameters.
    Average error 6.5\%.

  \item \textbf{Predictions Database.} 240 Seebeck predictions,
    filterable by element, cluster size, lattice type, viability.
    Sortable columns.

  \item \textbf{Account System.} Clerk authentication (Google OAuth),
    account page showing credits, usage, plan type.

  \item \textbf{Billing.} Stripe integration. Per-analysis (\$25),
    Academic (\$5K/yr), Industry (\$50K/yr). Credits tracked in
    DynamoDB. Failed analyses are not charged.

  \item \textbf{Security.}
    \begin{itemize}
      \item SCE-SIV encryption of all internal computation state
      \item T2 code tampering detection (12 files hashed at startup)
      \item 30+ trade secret terms blocked from LLM narrator
      \item All internal field names removed from API responses
      \item Rate limiting (10 burst, 5/s sustained)
    \end{itemize}

  \item \textbf{Multi-Domain Messaging.} Site addresses physics,
    pharma, battery, semiconductor, medical, and geological
    audiences. Pharma DSC example shown alongside physics examples.
\end{enumerate}

\subsection{Infrastructure}

\begin{itemize}
  \item Domain: dsf-ai.com (Porkbun, DNS to CloudFront)
  \item CDN: CloudFront distribution E17JT9XGBFU493
  \item SSL: ACM certificate (auto-renewing)
  \item API: API Gateway (HTTPS) $\to$ Lambda (cluster screener,
    billing, usage tracking) + ECS Fargate (CSV analysis)
  \item Auth: Clerk production instance, Google OAuth
  \item Billing: Stripe (sandbox; switch to live for real payments)
  \item Storage: DynamoDB (usage tracking), S3 (static site, SES
    envelopes)
  \item Email: support@dsf-ai.com (Porkbun forwarding)
\end{itemize}

\subsection{Pharma DSC Capability}

On May 11, 2026, the service was tested on simulated
pharmaceutical DSC (Differential Scanning Calorimetry) data.
All five transitions were detected:

\begin{center}
\begin{tabular}{@{}llll@{}}
\toprule
\textbf{Transition} & \textbf{Expected} & \textbf{Detected} &
  \textbf{Regime} \\
\midrule
Glass transition & 55$^\circ$C & 62$^\circ$C & VOLATILE \\
Polymorph I & 86$^\circ$C & 86$^\circ$C & VOLATILE \\
Polymorph II & 143$^\circ$C & 143$^\circ$C & VOLATILE \\
Melt & 211$^\circ$C & 211$^\circ$C & VOLATILE \\
Decomposition & 265$^\circ$C & 263$^\circ$C & VOLATILE \\
\bottomrule
\end{tabular}
\end{center}

This confirms DSF-AI works on DSC heat flow data, not just
resistance or dielectric measurements. Polymorph screening is
an FDA regulatory requirement, creating non-optional demand
in the pharmaceutical market.

\subsection{Data Normalization}

The service normalizes all input data before analysis:
$x_{\text{norm}} = ((x - x_{\min}) / (x_{\max} - x_{\min}))
\times 1000 + 1$. This ensures the analysis pipeline handles
any input range (milliohms, milliwatts, gigapascals, microvolts)
without domain-specific configuration.

\subsection{Remaining Work}

\begin{itemize}
  \item Sensor coupling weights tab (ArcLoom calibration data)
  \item Admin dashboard (usage statistics, customer management)
  \item Switch Stripe to live mode for real payments
  \item Stable ECS endpoint (ALB instead of ephemeral IP)
  \item PDF report download
\end{itemize}

\end{document}
